\chapter{Conclusão}
\label{chap:conclusão}

Esta monografia de proposta apresentou a contextualização, a fundamentação teórica e a metodologia para a avaliação de algoritmos de agregação robusta em Aprendizado Federado (\gls{fl}) sob cenários adversariais, com ênfase em ataques de envenenamento de modelo e no impacto da heterogeneidade estatística (\gls{noniid}). Ao longo do trabalho, foram delimitados o problema de pesquisa, os objetivos e o escopo experimental, destacando a relevância de mecanismos de agregação capazes de reduzir a influência de atualizações anômalas sem comprometer o desempenho em condições realistas de dados distribuídos. 

A fundamentação teórica consolidou conceitos centrais de \gls{fl}, vulnerabilidades na etapa de agregação e categorias de ataques, além de discutir estratégias de defesa por agregação robusta, incluindo métodos clássicos e abordagens modernas. Esse embasamento sustenta a necessidade de protocolos de avaliação mais rigorosos e comparáveis, uma vez que escolhas experimentais (como o nível de heterogeneidade, a fração de clientes maliciosos e o conjunto de ataques) podem influenciar significativamente a interpretação de robustez de um método.

No que se refere à metodologia proposta, foi planejado um ambiente de simulação reprodutível, com controle de \textit{seeds}, definição de um modelo local fixo e um protocolo experimental que permite comparar, sob as mesmas condições, um agregador de referência (FedAvg) e diferentes alternativas robustas (Krum, \textit{Trimmed Mean}, \textit{Coordinate-wise Median}, FLRAM, Mini-FL e MinVar). Também foram especificados os conjuntos de dados e o particionamento \gls{noniid} por Dirichlet, além da variação da fração de participantes maliciosos e dos ataques a serem implementados (por exemplo, \textit{sign-flip}, \textit{scaling}, \textit{MinMax} e um ataque adaptativo), permitindo analisar robustez, estabilidade do treinamento e custo computacional de forma sistemática.

Dessa forma, conclui-se que o planejamento apresentado é tecnicamente viável e fornece uma base consistente para a etapa seguinte do trabalho, que consiste na implementação do ambiente federado, dos agregadores e dos ataques, seguida da execução dos experimentos e das análises previstas. Espera-se que a continuidade do estudo resulte em evidências mais claras sobre limites de robustez e \textit{trade-offs} entre desempenho, heterogeneidade e adversarialidade, contribuindo para diretrizes práticas de seleção de agregadores em sistemas federados com participantes potencialmente não confiáveis.
